\section{Titling.sty and Hyperlinks.sty}\label{sec:ref-titling}

Two small base packages, both built on the same hook pattern: plain
defaults here, a class overrides with \cs{renewcommand}.

\subsection*{Titling}

The shared title block for LectureNotes and CourseDocument -- title,
subtitle, ``Updated $\langle$date$\rangle$'' -- structure written once,
class differences reduced to hooks. (Assessment deliberately does not load
it: an assessment opens with \cs{AssessmentTitle}, a different artefact.)
Also home of \file{datetime2}, the date's formatter.

\begin{center}
\begin{tabularx}{\linewidth}{@{}l Y@{}}
\cs{subtitle}\texttt{\{...\}} & the subtitle line(s); set before
  \cs{maketitle}.\\
\cs{TitleColor} & a \emph{style} hook (e.g.\
  \cs{color}\texttt{\{HeadingColor\}}), deliberately not a colour name --
  it keeps Titling free of any xcolor dependency. Default: empty.\\
\cs{TitleStamp} & follows the date (CourseDocument sets ``at
  $\langle$time$\rangle$''). Default: empty.\\
\cs{TitleMarker} & follows the stamp; currently a free hook -- the
  instructor-copy signal moved to the page frame. Default: empty.\\
\end{tabularx}
\end{center}

\subsection*{Hyperlinks}

hyperref plus a back-to-top anchor. The load itself is the point:
hyperref patches much of LaTeX, so anything loaded after it can silently
undo those patches -- Hyperlinks therefore defers it with
\cs{AtEndPreamble}, which runs after every package the \emph{document}
loads too, not merely after the class
(section~\ref{sec:load-order}). Requires the class to have passed
\texttt{dvipsnames} to xcolor beforehand.

\begin{center}
\begin{tabularx}{\linewidth}{@{}l Y@{}}
\cs{LinkColor} & internal links (default \texttt{black}; CourseDocument
  sets its identity colour). Read late, so overrides anywhere in the
  preamble win.\\
\cs{UrlColor} & external URLs (default \texttt{RoyalBlue}).\\
\cs{BackToTopLink}\texttt{\{...\}} & wrap anything in a jump to the top
  anchor (the projector footer's grey arrow).\\
\cs{TopAnchorName} & the anchor's name (\texttt{CourseTop}), placed at
  the very top of the body.\\
\cs{demo}\texttt{[tool]\{URL\}} & a link out to an interactive tool; see
  below.\\
\cs{DemoLabel}, \cs{DemoMarker} & the student-facing wording
  (\texttt{interactive graph}) and its cue ($\triangleright$).\\
\cs{DemoFormat}\texttt{\{URL\}} & the whole rendering, in one hook.\\
\end{tabularx}
\end{center}

\subsection*{\texorpdfstring{\cs{demo}}{\textbackslash demo}}

\emph{Students do not care which tool it is; only the author does.} A
class that meets ``LINK: DESMOS'', ``Desmos demo'' and ``LET'S EXPLORE
WITH DESMOS!'' on three consecutive pages is reading three different
things when it is always the same thing: press here, a graph moves. So
every demo prints \emph{one} uniform label:

\begin{manexsource}
\demo{https://www.desmos.com/calculator/xxxxxxxxxx}
\demo[geogebra]{https://example.org/some/opaque/url}
\end{manexsource}

The optional argument is \emph{accepted and discarded} -- the same
store-only idiom as ProblemMeta -- naming the tool for a human reading
the source, and for the link-checker when the domain is too opaque to
classify. Nothing about it reaches the page.

\textbf{Write the whole URL, scheme and all.} A URL without
\texttt{://} is a \emph{relative} one: the viewer resolves it against the
document's own location -- a local file path -- and the link silently does
nothing, while the page still looks perfectly right. Pasting
\texttt{www.example.com} without its \texttt{https://} is the entire failure,
so \cs{demo} checks and warns at build time rather than letting it through:

\begin{manexsource}
Package Hyperlinks Warning: \demo URL has no scheme and will
(Hyperlinks)                not be clickable: www.cbc.ca on input line 29.
\end{manexsource}

Same principle as TextRef's unresolved keys (section~\ref{sec:ref-textref}):
degrade if you must, but never silently.

The real payoff is enumerability: with every interactive link routed
through one macro, ``curl all my demos and tell me which died over the
summer'' is a grep away. That is reason enough to write \cs{demo} rather
than the bare \cs{href} you were about to.

It lives in Hyperlinks rather than TextRef (section~\ref{sec:ref-textref})
deliberately: a demo has nothing to do with the textbook -- it is as often
your own Desmos page -- and Hyperlinks is loaded by LectureNotes and
CourseDocument alike, so a syllabus gets \cs{demo} too. Assessment loads
no hyperref and therefore has none.
